% SEGMENT OLP-0066-S01
% Part: first-order-logic
% Chapter: proof-systems
% Section: natural-deduction

\documentclass[../../../include/open-logic-section]{subfiles}

\begin{document}

\iftag{FOL}
      {\olfileid{fol}{prf}{ntd}}
      {\olfileid{pl}{prf}{ntd}}

\olsection{இயல்பான வருவித்தல்}

இயல்பான வருவித்தல் என்பது நடைமுறைப் பகுத்தறிதலை (குறிப்பாகக் கணிதவியலார்
பயன்படுத்தும் ஒழுங்குபடுத்தப்பட்ட பகுத்தறிதலை) பிரதிபலிக்க நோக்கம்கொண்ட
!!a{derivation} முறைமை. நடைமுறைப் பகுத்தறிதல் பல ``இயல்பான''
பாங்குகளில் நடைபெறுகிறது. எடுத்துக்காட்டாக, ஒரு பிரிப்பிணைவு முன்கோளின்
ஒவ்வொரு பிரிப்புறுப்பிலிருந்தும் முடிவு பின்வருகிறது என்று நிறுவுவதன் மூலம்,
அந்த முன்கோளை அடிப்படையாகக் கொண்டு முடிவை நிறுவ நேர்வுவாரி நிறுவல்
அனுமதிக்கிறது. ஒரு முடிவின் மறுப்பு முரண்பாட்டிற்கு இட்டுச்செல்கிறது என்று
காட்டுவதன் மூலம், அந்த முடிவை நிறுவ மறைமுக நிறுவல் அனுமதிக்கிறது. முன்னிலையிலிருந்து
பின்னிலை பின்வருகிறது என்று காட்டுவதன் மூலம் ``எனில் \dots அப்போது \dots''
என்ற நிபந்தனைக் கூற்றை நிபந்தனை நிறுவல் நிறுவுகிறது. இத்தகைய இயல்பான
அனுமானங்களில் சிலவற்றின் முறையாக்கமே இயல்பான வருவித்தல். ஒவ்வொரு தருக்க
இணைப்பிற்கும் அளவையடைக்கும் இரண்டு விதிகள் உள்ளன: ஓர் அறிமுக விதியும்
ஒரு நீக்கல் விதியும். அவை ஒவ்வொன்றும் இத்தகைய இயல்பான அனுமானப் பாங்கு
ஒன்றுக்கு ஒத்திருக்கின்றன. எடுத்துக்காட்டாக, $\Intro{\lif}$ நிபந்தனை
நிறுவலுக்கும், $\Elim{\lor}$ நேர்வுவாரி நிறுவலுக்கும் ஒத்திருக்கின்றன.
$!A \land !B$ இலிருந்து~$!A$ ஐ (அல்லது $!B$ ஐ) அனுமானிக்க அனுமதிக்கும்
$\Elim{\land}$ குறிப்பாக எளிய விதி.

% SEGMENT OLP-0066-S02
எடுகோள்களைப் பயன்படுத்துவது இயல்பான வருவித்தலை மற்ற !!{derivation}
முறைமைகளிலிருந்து வேறுபடுத்தும் ஒரு பண்பு. இயல்பான வருவித்தலில்
!!^a{derivation} என்பது !!{formula}s இன் மரம். !!{formula}s இன் மரத்தின்
வேரில் ஒரு !!{formula} இருக்கும்; மரத்தின் ``இலைகள்'', முடிவு
வருவிக்கப்படும் !!{formula}s ஆகும். இயல்பான வருவித்தலில் சில இலை
!!{formula}s, !!{derivation}வினுள் பங்காற்றினாலும், !!{derivation} முடிவை
அடையும் நேரத்திற்குள் அவற்றின் ``பயன் தீர்ந்துவிடும்.'' குறுகிய நேரத்திற்கு
மட்டுமே செயலிலிருக்கும் கருதுகோள்களை நடைமுறைப் பகுத்தறிதலில் அறிமுகப்படுத்தும்
வழக்கை இது பிரதிபலிக்கிறது. எடுத்துக்காட்டாக, நேர்வுவாரி நிறுவலில் ஒவ்வொரு
பிரிப்புறுப்பும் மெய் என்று எடுகொள்கிறோம்; நிபந்தனை நிறுவலில் முன்னிலை மெய்
என்று எடுகொள்கிறோம்; மறைமுக நிறுவலில் முடிவின் மறுப்பு மெய் என்று
எடுகொள்கிறோம். ஓர் இடைநிலைப் படியை நிறுவுவதற்காகக் கருதுகோள் எடுகோள்களை
அறிமுகப்படுத்திப் பின்னர் அவற்றை நீக்கும் இம்முறை இயல்பான வருவித்தலின்
தனிச்சிறப்பு. இயல்பான வருவித்தல் !!{derivation}வின் இலைகளிலுள்ள
வாய்பாடுகள் எடுகோள்கள் எனப்படுகின்றன; சில அனுமான விதிகளில் அவற்றின்
``!!{discharge}'' நிகழலாம். எடுத்துக்காட்டாக,~$!A$ ஐ உள்ளடக்கிய சில
எடுகோள்களிலிருந்து~$!B$ க்கு !!a{derivation} நம்மிடம் இருந்தால்,
$!A \lif !B$ ஐ அனுமானிக்கவும் $!A$ வடிவிலுள்ள எந்த எடுகோளையும்
விடுவிக்கவும் $\Intro{\lif}$ விதி அனுமதிக்கிறது. (எந்த அனுமானங்களில் எந்த
எடுகோள்கள் விடுவிக்கப்படுகின்றன என்பதைக் கண்காணிக்க, அனுமானத்திற்கும் அது
விடுவிக்கும் எடுகோள்களுக்கும் ஓர் எண்ணால் பெயரிடுகிறோம்.) !!{derivation}வின்
முடிவில் !!{undischarged} எடுகோள்கள் அனைத்தும் சேர்ந்து முடிவின் மெய்மைக்குப்
போதுமானவை. ஆகவே !!a{derivation}, அதன் !!{undischarged} எடுகோள்கள் அதன்
முடிவைப் பொருண்மைப்படிப் பின்விளைவிக்கின்றன என்பதை நிறுவுகிறது.

% SEGMENT OLP-0066-S03
இயல்பான வருவித்தலை அடிப்படையாகக் கொண்ட $\Gamma \Proves !A$ உறவு
பொருந்துவது, அப்போதும் அப்போது மட்டுமே, $!A$~மரத்தின் கடைசி
!!{sentence} ஆகவும், !!{undischarged} ஒவ்வொரு இலையும்~$\Gamma$ இல்
உள்ளதாகவும் இருக்கும் !!a{derivation} ஒன்று உண்டு. இயல்பான வருவித்தலில்
$!A$~ஒரு தேற்றம் என்பது, அப்போதும் அப்போது மட்டுமே, $!A$~கடைசி
!!{sentence} ஆகவும் எல்லா எடுகோள்களும் !!{discharged} ஆகவும் இருக்கும்
!!a{derivation} ஒன்று உண்டு. எடுத்துக்காட்டாக, $\Proves (!A
\land !B) \lif !A$ என்பதைக் காட்டும் !!a{derivation} இதோ:
\begin{prooftree}
  \AxiomC{$\Discharge{!A \land !B}{1}$}
  \RightLabel{\Elim{\land}}
  \UnaryInfC{$!A$}
  \DischargeRule{\Intro{\lif}}{1}
  \UnaryInfC{$(!A \land !B) \lif !A$}
\end{prooftree}
$!A \land !B$ என்ற எடுகோள் \Intro{\lif} அனுமானத்தில்
!!{discharged} என்பதைப் பெயர்~$1$ குறிக்கிறது.

% SEGMENT OLP-0066-S04
கணம்~$\Gamma$ இயல்பான வருவித்தலில் முரணுடையது என்பது, அப்போதும் அப்போது
மட்டுமே, $\Gamma \Proves \lfalse$. முரணுடைய கணத்திலிருந்து எந்த
!!{sentence}வையும் !!{derive} இயலுமாறு \FalseInt{} விதி செய்கிறது.

% SEGMENT OLP-0066-S05
இயல்பான வருவித்தல் முறைமைகளை 1930களில் கெர்ஹார்ட் கென்ட்சன் மற்றும்
% RENDER LOCALIZATION TA-RENDER-001: transliterate the Polish name in Tamil.
ஸ்டானிஸ்லாவ் யாஸ்கோவ்ஸ்கி உருவாக்கினர்; பின்னர் டாக் ப்ராவிட்ஸ் மற்றும்
ஃப்ரெடெரிக் ஃபிட்ச் அவற்றை மேலும் வளர்த்தனர். அவற்றின் அனுமானங்கள் இயல்பான
நிறுவல் முறைகளைப் பிரதிபலிப்பதால், மெய்யியலார் அவற்றை விரும்புகின்றனர்.
ஃபிட்ச் உருவாக்கிய வடிவங்கள் அறிமுகத் தருக்கவியல் பாடநூல்களில் அடிக்கடி
பயன்படுகின்றன. தருக்க மெய்யியலில், இயல்பான வருவித்தலின் விதிகளே தருக்கச்
செயலிகளின் பொருள்களைத் தருவதாகச் சிலவேளைகளில் கொள்ளப்பட்டுள்ளன
(``நிறுவல் கோட்பாட்டுப் பொருண்மையியல்'').

\end{document}
